% Application letter - Humboldt-Universitaet zu Berlin, Research Fellow (PhD), ref. DR/153/26
% "Hybrid-plasmonic platform with tuneable light-matter coupling in a 2D parameter space"
% Compile with pdfLaTeX.
\documentclass[11pt,a4paper]{article}
\usepackage[utf8]{inputenc}
% T1 + Latin Modern when available; otherwise default Computer Modern.
\IfFileExists{lmodern.sty}{\usepackage[T1]{fontenc}\usepackage{lmodern}}{}
\usepackage[top=1.4cm,bottom=1.3cm,left=2cm,right=2cm]{geometry}
\usepackage{xcolor}
\usepackage{microtype}
\usepackage[hidelinks]{hyperref}

\pagestyle{empty}
\setlength{\parindent}{0pt}
\setlength{\parskip}{5pt}
\definecolor{muted}{gray}{0.38}
\newcommand{\code}[2]{\href{#1}{\underline{#2}}}

\begin{document}

% ---------- Header (same as the CV) ----------
{\LARGE\bfseries Matteo Falasca}\par\vspace{-2pt}
{\small\color{muted}MSc student in theoretical physics, University of Trieste\par\vspace{-5pt}}
\vspace{-6pt}\noindent\rule{\linewidth}{0.3pt}\par

\vspace{3pt}
{\small
\begin{minipage}[t]{0.68\linewidth}
PD Dr.\ Sylke Blumstengel\\
Department of Physics, Humboldt-Universität zu Berlin
\end{minipage}\hfill
\begin{minipage}[t]{0.3\linewidth}\raggedleft\color{muted}
\today
\end{minipage}}

\vspace{5pt}
\textbf{PhD application ``Hybrid-plasmonic platform with tuneable light--matter coupling in a 2D parameter space'' job ID DR/153/26 }

Dear Dr.\ Blumstengel,

I am completing my MSc in theoretical physics at the University of Trieste,
where I expect to graduate in December 2026, and I am applying
for the doctoral position in your group. My training is in
light--matter interactions, many-body physics and
numerical computation.
I am applying for this position because I want
to continue my studies in a field
where theory can be simulated and where my knowledge
can be useful to real
state-of-the-art experiments and implementations.

Light--matter interaction is at the centre of my
Master's thesis, ``Graviton Scattering and Induced Decoherence'',
supervised by Prof.\ Angelo Bassi. I study how a bosonic field that
scatters off matter changes the dynamics of the matter itself. I
quantized the field coupled to matter and derived the scattering
cross sections, which generalize the light--matter case and include
velocity corrections I could not find in the literature. The work
gave me a solid command of quantum field theory and second
quantization, and a clear picture of how the coupling to a
field defines the behaviour of matter: the same physics
your platform makes tuneable.

Your project aims to demonstrate
the potential of plasmonic platforms, partly
through numerical finite element simulations,
and simulation and data analysis are where my experience is
strongest. In my Bachelor's thesis I
analysed eROSITA data and built a
neural-network emulator for a parameter
space too large to sample directly, using
Latin hypercube sampling.
Sampling a parameter
space efficiently would be a huge resource for
a platform tuneable
in a 2D parameter space like yours.
For the noisy-circuit work of Grossi
et al.\ (arXiv:2410.01893)
I wrote a faster algorithm that computes information
flows and that was
later used in their research,
and in a computational physics course
I acquired experience from Quantum Monte Carlo  
to simulating systems with
emerging collective behaviours.
I code in Python and Fortran and can use software
tools for data processing and analysis.

I built my background in condensed matter physics
and optics through my choice of courses, starting with
Atoms, Molecules and Photons with Prof.\ Francesco Scazza
and Quantum Many-Body Systems with Prof.\ Federico Becca,
neither of which was mandatory for my curriculum, together with
several courses on Statistical Mechanics and Quantum Field Theory.
My interest in these topics grew in
2025, when I was selected for the three-week ICTP School on
Quantum Dynamics of Matter, Light and Information, and when
I attended the DPG Bad Honnef school on quantum metrology,
both closely related with light--matter interactions.

I have not yet used optical spectroscopy in the lab,
but I have met it on the theory side. Among my study projects,
the one for Atoms, Molecules and Photons consisted of
presenting research results at a domestic conference on
the creation of noise-resistant entangled
states of light by the addition of a single delocalized photon, and it
led me to study homodyne detection
and how quantum states of
light are measured.
This is a part of physics I would like to
get into: performing measurements and building
theory next to verifiable experiments fascinates me,
 I am looking for the chance to do it.
Personally, I would feel very comfortable presenting
results in international journals and at conferences:
I have been used to public speaking since I was young, 
and I stay calm in front of an audience.

On a personal level, I am an outgoing person 
and I like working with others, with
whom I build healthy working environments.
I tutored first-year Bachelor students,
and I have refereed football since 2017,
now in the top regional division, which taught me
to decide under pressure, to work in a team and to 
explain my choices.

Prof.\ Angelo Bassi and Prof.\ Francesco Scazza have agreed to write reference letters for me and can send them directly to you. Thank you for your time and consideration.

Yours sincerely,

Matteo Falasca \hfill {\small\color{muted}Code:\quad
\code{https://github.com/MatteoFalasca-AFK/Projects.CV/tree/main}{Some of my projects}\quad
\end{document}